\chapter{رشته‌ها، زبان‌ها و روش‌های اثبات}

همه ماشین‌هایی که در این درس می‌بینیم ورودی خود را به شکل یک \textbf{رشته} می‌گیرند. مجموعه‌ای از رشته‌ها را \textbf{زبان} می‌نامیم. در نیمه اول این فصل این دو مفهوم را دقیق تعریف می‌کنیم. در نیمه دوم روش‌های اصلی اثبات را مرور می‌کنیم، چون بخش بزرگی از این درس ثابت کردن قضیه‌ها درباره زبان‌ها و ماشین‌هاست.

\section{الفبا، رشته و زبان}

\begin{definition}[الفبا و رشته]
\begin{itemize}
\item \term{الفبا}{Alphabet} هر مجموعه \textbf{متناهی و ناتهی} است. الفبا را معمولاً با $\Sigma$ یا $\Gamma$ نشان می‌دهیم و به عضوهای آن \textbf{نماد} می‌گوییم.
\item \term{رشته}{String} روی یک الفبا، دنباله‌ای \textbf{متناهی} از نمادهای آن الفباست. نمادها را بدون فاصله و ویرگول پشت هم می‌نویسیم.
\end{itemize}
\end{definition}

\begin{example}
$\Sigma_1=\{0,1\}$، $\Sigma_2=\{a,b,c,\dots,z\}$ و $\Gamma=\{0,1,x,y,z\}$ سه الفبا هستند. رشته $01001$ روی $\Sigma_1$ و رشته $abracadabra$ روی $\Sigma_2$ است.
\end{example}

چند مفهوم پایه درباره رشته‌ها:
\begin{itemize}
\item \textbf{طول:} اگر $w=w_1w_2\cdots w_n$ باشد ($w_i\in\Sigma$)، طول آن $\abs{w}=n$ است.
\item \textbf{رشته تهی} $\eps$ رشته‌ای با طول صفر است. (رشته تهی همان نقش عدد صفر در جمع را دارد.)
\item \textbf{معکوس:} $w^R=w_nw_{n-1}\cdots w_1$. مثلاً $abb^R=bba$، $a^R=a$ و $\eps^R=\eps$.
\item \textbf{زیررشته:} $z$ زیررشته $w$ است اگر $z$ به صورت \emph{پیوسته} درون $w$ بیاید. مثلاً $cad$ زیررشته $abracadabra$ است، $0100$ زیررشته $001000$ است، اما $0000$ زیررشته $001000$ نیست.
\item \textbf{الحاق:} اگر $x=x_1\cdots x_m$ و $y=y_1\cdots y_n$، آنگاه $xy=x_1\cdots x_my_1\cdots y_n$. داریم $\abs{xy}=\abs{x}+\abs{y}$ و $\eps x=x\eps=x$. همچنین $x^k=\underbrace{xx\cdots x}_{k}$، $x^0=\eps$ و $\abs{x^k}=k\abs{x}$.
\item \textbf{پیشوند:} $x$ پیشوند $y$ است اگر رشته‌ای مثل $z$ باشد که $xz=y$. اگر علاوه بر این $x\neq y$، به $x$ \textbf{پیشوند سره} می‌گوییم. \textbf{پسوند} هم به همین شکل تعریف می‌شود: $x$ پسوند $y$ است اگر $zx=y$.
\end{itemize}

\begin{example}
پیشوندهای $aaba$ عبارت‌اند از $\eps$، $a$، $aa$، $aab$ و $aaba$. به طور کلی رشته‌ای به طول $n$ دقیقاً $n+1$ پیشوند دارد (پیشوندهایی به طول $0,1,\dots,n$). تنها پیشوند $\eps$ خودش است.
\end{example}

\subsection{ترتیب رشته‌ها}
\textbf{ترتیب لغت‌نامه‌ای} همان ترتیب واژه‌ها در فرهنگ لغت است. اما در این درس معمولاً از \term{ترتیب رشته‌ای}{Shortlex Order} استفاده می‌کنیم: رشته‌های کوتاه‌تر جلوتر می‌آیند و رشته‌های هم‌طول به ترتیب لغت‌نامه‌ای مرتب می‌شوند:
\[(\eps,\ 0,\ 1,\ 00,\ 01,\ 10,\ 11,\ 000,\ \dots).\]

\begin{note}
چرا ترتیب لغت‌نامه‌ای کافی نیست؟ در ترتیب لغت‌نامه‌ای روی $\{0,1\}$ فهرست با $\eps,0,00,000,\dots$ شروع می‌شود و هیچ‌وقت به رشته $1$ نمی‌رسد. اما در ترتیب رشته‌ای هر رشته بالاخره در یک جای مشخص فهرست قرار می‌گیرد. از این ویژگی بعداً برای فهرست کردن همه رشته‌ها استفاده می‌کنیم (مثلاً برای اثبات شمارا بودن $\Sigma^*$).
\end{note}

\subsection{زبان‌ها}

\begin{definition}[زبان]
\term{زبان}{Language} هر مجموعه‌ای از رشته‌هاست. $\Sigma^*$ مجموعه همه رشته‌ها روی $\Sigma$ (از جمله $\eps$) و $\Sigma^+=\Sigma^*-\{\eps\}$ است. پس هر زبان روی $\Sigma$ زیرمجموعه‌ای از $\Sigma^*$ است. زبانی که هیچ عضوش پیشوند سره عضو دیگری نباشد \textbf{بی‌پیشوند} (\lr{Prefix-free}) نام دارد.
\end{definition}

\begin{example}
زبان‌های متناهی: $L_1=\emptyset$، $L_2=\{\eps\}$ و $L_3=\{a,aa,aba\}$. زبان‌های نامتناهی: $L_4=\Sigma^*$، $L_5=\Sigma^+$ و
\[L_6=\{a^nb^n\st n\ge0\}=\{\eps,ab,aabb,aaabbb,\dots\}.\]
\end{example}

\begin{warning}
$\emptyset$ و $\{\eps\}$ یکی نیستند! $\emptyset$ هیچ رشته‌ای ندارد، اما $\{\eps\}$ یک عضو دارد: رشته تهی.
\end{warning}

\begin{definition}[عملیات روی زبان‌ها]
فرض کنید $L,L_1,L_2$ زبان‌هایی روی $\Sigma$ باشند.
\begin{itemize}
\item \textbf{اجتماع و اشتراک:} چون زبان‌ها مجموعه‌اند، همان اجتماع و اشتراک مجموعه‌ها.
\item \textbf{متمم:} $\overline{L}=\Sigma^*-L$.
\item \textbf{معکوس:} $L^R=\{w^R\st w\in L\}$.
\item \textbf{الحاق:} $L_1\circ L_2=\{xy\st x\in L_1,\ y\in L_2\}$.
\item \textbf{ستاره کلینی:} $L^*=\{x_1x_2\cdots x_k\st k\ge0,\ x_i\in L\}$. چون $k=0$ مجاز است، همیشه $\eps\in L^*$.
\end{itemize}
\end{definition}

جدول زیر سه مفهوم اصلی را مقایسه می‌کند:
\begin{center}
\begin{tabular}{lcc}
\toprule
 & می‌تواند تهی باشد؟ & می‌تواند نامتناهی باشد؟\\
\midrule
الفبا & \textcolor{warnFrame}{$\times$} & \textcolor{warnFrame}{$\times$}\\
رشته & \textcolor{exFrame}{$\checkmark$} & \textcolor{warnFrame}{$\times$}\\
زبان & \textcolor{exFrame}{$\checkmark$} & \textcolor{exFrame}{$\checkmark$}\\
\bottomrule
\end{tabular}
\end{center}

\section{تعریف، قضیه و اثبات}

\begin{definition}[واژه‌های پایه ریاضی]
\begin{itemize}
\item \textbf{تعریف} شیئی را که به کار می‌بریم و نماد آن را توضیح می‌دهد. تعریف باید دقیق باشد، چه برای مفهومی ساده مثل «مجموعه» و چه برای مفهومی پیچیده مثل «امنیت».
\item \textbf{گزاره ریاضی} ویژگی‌ای از یک شیء را بیان می‌کند که ممکن است درست یا نادرست باشد، اما نباید مبهم باشد.
\item \textbf{اثبات} استدلالی منطقی و قانع‌کننده است که درستی یک گزاره را بدون هیچ شکی نشان می‌دهد.
\item \textbf{قضیه} گزاره‌ای است که درستی آن اثبات شده است. قضیه‌های کمکی را \textbf{لم} و قضیه‌هایی را که به‌سادگی از قضیه دیگری نتیجه می‌شوند \textbf{نتیجه} می‌نامیم.
\end{itemize}
\end{definition}

\subsection{چطور اثبات پیدا کنیم؟}
پیدا کردن اثبات همیشه ساده نیست، اما این راهکارها کمک می‌کنند:
\begin{enumerate}
\item نمادها را کامل بفهمید و گزاره را با کلمات خودتان بازنویسی کنید.
\item گزاره را به بخش‌های کوچک‌تر بشکنید:
  \begin{itemize}
  \item «$P$ اگر و تنها اگر $Q$» یعنی $P\leftrightarrow Q$ و دو جهت دارد: جهت رفت $P\rightarrow Q$ و جهت برگشت $Q\rightarrow P$.
  \item «مجموعه‌های $A$ و $B$ برابرند» یعنی هر عضو $A$ در $B$ است و هر عضو $B$ در $A$ است.
  \end{itemize}
\item دنبال مثال نقض بگردید. اگر پیدا نشد، معمولاً دلیلش شما را به اثبات می‌رساند.
\item حالت‌های ساده‌تر را امتحان کنید.
\end{enumerate}

\begin{example}
گزاره: «مجموع درجه‌های همه رأس‌های یک گراف بی‌جهت عددی زوج است.» با امتحان چند گراف کوچک (بدون یال: مجموع $0$؛ یک یال: $2$؛ دو یال: $4$؛ سه یال: $6$) می‌بینیم که \textbf{هر بار یک یال اضافه می‌کنیم، مجموع دقیقاً دو واحد زیاد می‌شود}. همین مشاهده ما را به اثبات (قضیه \ref{thm:degree}) می‌رساند.
\end{example}

\subsection{چطور اثبات بنویسیم؟}
یک اثبات خوب دنباله‌ای از جمله‌هاست که هر کدام به‌طور منطقی از قبلی‌ها نتیجه می‌شود. دقیق باشید، مرتب بنویسید و از شکل و متن ساده کمک بگیرید. اگر اثبات طولانی است، اول یک طرح کلی بدهید و بعد جزئیات را بنویسید. خلاصه: \textbf{مرتب و مختصر}.

\begin{theorem}[دمورگان]
برای هر دو مجموعه $A$ و $B$ داریم $\overline{A\cup B}=\overline{A}\cap\overline{B}$.
\end{theorem}
\begin{proof}
برابری دو مجموعه را با اثبات دو جهت نشان می‌دهیم.

\textbf{جهت رفت} ($\overline{A\cup B}\subseteq\overline{A}\cap\overline{B}$): فرض کنید $x\in\overline{A\cup B}$. پس $x\notin A\cup B$؛ یعنی $x$ نه در $A$ است و نه در $B$. پس $x\in\overline{A}$ و $x\in\overline{B}$، یعنی $x\in\overline{A}\cap\overline{B}$.

\textbf{جهت برگشت} ($\overline{A}\cap\overline{B}\subseteq\overline{A\cup B}$): فرض کنید $x\in\overline{A}\cap\overline{B}$. پس $x\notin A$ و $x\notin B$. بنابراین $x$ در اجتماع آن‌ها هم نیست: $x\notin A\cup B$، یعنی $x\in\overline{A\cup B}$.
\end{proof}

\begin{theorem}\label{thm:degree}
در هر گراف بی‌جهت $G$ مجموع درجه‌های همه رأس‌ها عددی زوج است.
\end{theorem}
\begin{proof}
فرض کنید $G=(V,E)$ و $d(v)$ درجه رأس $v$ باشد. هر یال $\{u,v\}\in E$ دقیقاً یک واحد به $d(u)$ و یک واحد به $d(v)$ اضافه می‌کند. پس هر یال دقیقاً $2$ واحد در مجموع درجه‌ها سهم دارد و
\[\sum_{v\in V}d(v)=\sum_{\{u,v\}\in E}(1+1)=2\abs{E},\]
که عددی زوج است.
\end{proof}

\section{انواع روش‌های اثبات}

یک اثبات ممکن است از چند روش مختلف در بخش‌های مختلفش استفاده کند. روش‌های اصلی این‌ها هستند:

\subsection{اثبات مستقیم و غیرمستقیم}
\textbf{اثبات مستقیم} بر قاعده‌ای به نام \lr{modus ponens} («روش تأیید») تکیه دارد: اگر $p$ و $p\rightarrow q$ ثابت شده باشند، $q$ هم ثابت شده است.
\textbf{اثبات غیرمستقیم} بر قاعده \lr{modus tollens} («روش نفی») تکیه دارد: اگر $\neg q$ و $p\rightarrow q$ ثابت شده باشند، $\neg p$ هم ثابت شده است.
\[
\begin{array}{l} p\rightarrow q\\ p\\ \hline \therefore\ q\end{array}
\qquad\qquad
\begin{array}{l} p\rightarrow q\\ \neg q\\ \hline \therefore\ \neg p\end{array}
\]
یک کاربرد رایج روش غیرمستقیم \textbf{عکس نقیض} است: به جای اثبات $p\rightarrow q$، گزاره هم‌ارز $\neg q\rightarrow\neg p$ را ثابت می‌کنیم. مثلاً برای اثبات «اگر $n^2$ زوج باشد، $n$ زوج است» کافی است نشان دهیم «اگر $n$ فرد باشد، $n^2$ فرد است»: اگر $n=2k+1$ آنگاه $n^2=2(2k^2+2k)+1$ فرد است.

\subsection{اثبات ساختنی}
وقتی می‌خواهیم وجود شیئی را ثابت کنیم ($\exists x\ P(x)$)، می‌توانیم نشان دهیم آن شیء را چطور بسازیم: یعنی $a$ را بسازیم و ثابت کنیم $P(a)$ درست است. در این درس این روش بسیار پرکاربرد است: برای اثبات منظم بودن یک زبان، ماشینی می‌سازیم که آن را بپذیرد.

\begin{theorem}
برای هر عدد زوج $n>2$ یک گراف $3$-منتظم (گرافی که درجه همه رأس‌هایش ۳ است) با $n$ رأس وجود دارد.
\end{theorem}
\begin{proof}
گراف $G=(V,E)$ را با $V=\{0,1,\dots,n-1\}$ و یال‌های زیر می‌سازیم:
\[
E=\underbrace{\{\{i,i+1\}\st 0\le i\le n-2\}\cup\{\{n-1,0\}\}}_{\text{\rl{یک دور}}}\ \cup\ \underbrace{\{\{i,i+n/2\}\st 0\le i\le n/2-1\}}_{\text{\rl{قطرها}}}.
\]
دسته اول یال‌ها رأس‌ها را روی یک دور قرار می‌دهد، پس هر رأس از این دسته دقیقاً ۲ یال دارد. دسته دوم هر رأس را به رأس روبه‌رویش روی دور وصل می‌کند. چون $n$ زوج است، $n/2$ عدد صحیح است و هر رأس دقیقاً یک رأس روبه‌رو دارد. دقیق‌تر: رأس $i<n/2$ با یال $\{i,i+n/2\}$ و رأس $j\ge n/2$ با یال $\{j-n/2,j\}$ وصل می‌شود، پس هر رأس دقیقاً یک قطر دارد. چون $n>2$، داریم $n/2\ge2$؛ پس رأس روبه‌رو هیچ‌وقت همسایه رأس روی دور نیست و یال تکراری ساخته نمی‌شود. در نتیجه درجه همه رأس‌ها $2+1=3$ است.
\end{proof}

\subsection{برهان خلف}
در برهان خلف فرض می‌کنیم گزاره نادرست است و نشان می‌دهیم این فرض به نتیجه‌ای نادرست (\textbf{تناقض}) می‌رسد. پس فرض باید غلط باشد. مثال روزمره: علی می‌بیند مریم کاملاً خشک از بیرون آمده است. اگر بیرون باران می‌بارید (فرض نادرست)، مریم خیس بود (نتیجه نادرست). پس باران نمی‌بارد.

\begin{theorem}
$\sqrt{2}$ گنگ است.
\end{theorem}
\begin{proof}
با برهان خلف فرض کنید $\sqrt2$ گویا باشد. پس $\sqrt2=m/n$ برای اعداد صحیح $m,n$. می‌توانیم کسر را ساده کنیم تا $m$ و $n$ نسبت به هم اول باشند؛ در نتیجه دست‌کم یکی از آن‌ها فرد است. با ضرب دو طرف در $n$ و به توان دو رساندن داریم
\[n\sqrt2=m\quad\Longrightarrow\quad 2n^2=m^2.\]
پس $m^2$ زوج است و در نتیجه $m$ زوج است (مثال عکس نقیض بالا). بنویسید $m=2k$. با جایگذاری: $2n^2=4k^2$، یعنی $n^2=2k^2$. پس $n$ هم زوج است. اما $m$ و $n$ هر دو زوج شدند، در حالی که فرض کرده بودیم دست‌کم یکی فرد است. تناقض.
\end{proof}

\subsection{استقرا}
استقرا روشی برای اثبات یک ویژگی برای همه عضوهای یک مجموعه نامتناهی مثل $\mathbb{N}$ است. اثبات دو قدم دارد:
\begin{itemize}
\item \textbf{پایه:} ویژگی را برای کوچک‌ترین عضو ثابت می‌کنیم، مثلاً $P(1)$. (لازم نیست حتماً از ۱ شروع کنیم.)
\item \textbf{گام استقرا:} فرض می‌کنیم $P(k)$ درست است (\textbf{فرض استقرا}) و ثابت می‌کنیم $P(k+1)$ هم درست است.
\end{itemize}
در \textbf{استقرای قوی} فرض استقرا این است که ویژگی برای همه مقادیر کوچک‌تر درست است: $\bigwedge_{i\in\{1,\dots,k\}}P(i)\rightarrow P(k+1)$.

\begin{example}[فرمول اقساط وام]
فرض کنید $P$ مبلغ وام، $P_t$ بدهی باقی‌مانده بعد از ماه $t$ام ($P_0=P$)، $I$ نرخ سود سالانه، $M=1+I/12$ ضریب سود ماهانه و $Y$ قسط ماهانه باشد. هر ماه بدهی به خاطر سود در $M$ ضرب و به اندازه قسط کم می‌شود: $P_t=MP_{t-1}-Y$.
\end{example}

\begin{theorem}
برای هر $t\ge0$:\quad $P_t=PM^t-Y\left(\dfrac{M^t-1}{M-1}\right)$.
\end{theorem}
\begin{proof}
با استقرا روی $t$.
\textbf{پایه} ($t=0$): $PM^0-Y\frac{M^0-1}{M-1}=P-0=P=P_0$.
\textbf{گام استقرا:} فرض کنید فرمول برای $k$ درست باشد. آنگاه
\begin{align*}
P_{k+1}=P_kM-Y
&=\left[PM^k-Y\left(\frac{M^k-1}{M-1}\right)\right]M-Y\\
&=PM^{k+1}-Y\left(\frac{M^{k+1}-M}{M-1}\right)-Y\left(\frac{M-1}{M-1}\right)\\
&=PM^{k+1}-Y\left(\frac{M^{k+1}-1}{M-1}\right).
\end{align*}
\end{proof}

\begin{theorem}
یک درخت دودویی به ارتفاع $n$ حداکثر $2^n$ برگ دارد.
\end{theorem}
\begin{proof}
فرض کنید $l(n)$ بیشترین تعداد برگ در یک درخت دودویی به ارتفاع $n$ باشد. (ارتفاع، طول بلندترین مسیر از ریشه تا یک برگ است.)
\textbf{پایه:} درختی به ارتفاع $0$ فقط ریشه را دارد، پس $l(0)=1\le2^0$.
\textbf{گام استقرا:} فرض کنید $l(j)\le 2^j$ برای هر $j\le k$. ریشه یک درخت به ارتفاع $k+1$ حداکثر دو فرزند دارد و زیردرخت هر فرزند ارتفاعی حداکثر $k$ دارد. برگ‌های درخت همان برگ‌های این زیردرخت‌ها هستند، پس
\[l(k+1)\le 2\cdot 2^k=2^{k+1}.\]
\end{proof}

\begin{theorem}
اگر $n$ خط راست در صفحه داشته باشیم که هر دو خط یکدیگر را قطع کنند و هیچ سه خطی از یک نقطه نگذرند، صفحه به
\[A(n)=\frac{n(n+1)}{2}+1\]
ناحیه تقسیم می‌شود.
\end{theorem}
\begin{proof}
\textbf{پایه:} یک خط صفحه را به $A(1)=2$ ناحیه تقسیم می‌کند و فرمول هم $\frac{1\cdot2}{2}+1=2$ می‌دهد.

\textbf{مشاهده کلیدی:} خط $(k+1)$ام همه $k$ خط قبلی را در $k$ نقطه متمایز قطع می‌کند. این $k$ نقطه، خط جدید را به $k+1$ تکه تقسیم می‌کنند و هر تکه یک ناحیه قبلی را به دو ناحیه می‌شکند. پس $A(k+1)=A(k)+k+1$.

\textbf{گام استقرا:} با فرض $A(k)=\frac{k(k+1)}{2}+1$:
\[A(k+1)=\frac{k(k+1)}{2}+1+k+1=\frac{k(k+1)+2(k+1)}{2}+1=\frac{(k+1)(k+2)}{2}+1.\]
\end{proof}

\begin{summary}
\begin{itemize}
\item الفبا مجموعه‌ای متناهی و ناتهی است؛ رشته دنباله‌ای متناهی از نمادهاست؛ زبان مجموعه‌ای (شاید نامتناهی) از رشته‌هاست.
\item عملیات اصلی روی زبان‌ها: اجتماع، اشتراک، متمم، معکوس، الحاق و ستاره.
\item ترتیب رشته‌ای (اول کوتاه‌ترها) اجازه می‌دهد همه رشته‌ها را فهرست کنیم.
\item روش‌های اثبات: مستقیم، غیرمستقیم (عکس نقیض)، ساختنی، برهان خلف و استقرا.
\end{itemize}
\end{summary}
